\section{Background}
\label{sec:background}

\subsection{Quantum Computing Concepts for QIEA}
\label{sec:background:quantum}

Quantum-inspired evolutionary algorithms borrow their representation and
update rules from the mathematics of a single qubit, without requiring
quantum hardware to run. A qubit's state is a normalized superposition of
the two classical basis states,
\begin{equation}
  \label{eq:qubit}
  \lvert\psi\rangle = \alpha\lvert 0\rangle + \beta\lvert 1\rangle,
  \qquad \lvert\alpha\rvert^2 + \lvert\beta\rvert^2 = 1,
\end{equation}
where $\lvert\alpha\rvert^2$ and $\lvert\beta\rvert^2$ are interpreted as
the probability of the qubit collapsing to $0$ or $1$ upon observation.
\citet{han2002qea} introduced the \emph{Q-bit chromosome}: an individual
is encoded not as a fixed binary string but as a string of such
probability-amplitude pairs $(\alpha_i,\beta_i)$, one per bit position.
Observing (measuring) a Q-bit chromosome collapses it into an ordinary
binary solution that can be evaluated by the problem's fitness function,
while the amplitudes themselves are what the algorithm evolves --
maintaining, in effect, a probabilistic superposition over the whole
population's search history rather than a single discrete solution.

The corresponding variation operator is the \emph{Q-gate}, a rotation
applied to each Q-bit's amplitude pair,
\begin{equation}
  \label{eq:qgate}
  \begin{pmatrix}\alpha_i'\\ \beta_i'\end{pmatrix}
  =
  \begin{pmatrix}\cos\theta_i & -\sin\theta_i\\ \sin\theta_i & \cos\theta_i\end{pmatrix}
  \begin{pmatrix}\alpha_i\\ \beta_i\end{pmatrix},
\end{equation}
where the rotation angle $\theta_i$ and its sign are looked up from a
predetermined table that compares the current bit's observed value and
fitness contribution against the best solution found so far, nudging the
amplitude pair toward whichever basis state that comparison favors. This
rotation-gate mechanism is what most surveyed QIEA methods share
regardless of application domain~\citep{vivek2025fss}, and it is the
representation/operator combination
Section~\ref{sec:taxonomy:representation} and
Section~\ref{sec:taxonomy:operator} use as the taxonomy's reference
point; variants that replace or augment it with interference- or
entanglement-inspired operators are discussed there rather than here.

\subsection{Multi-Objective Optimization Fundamentals}
\label{sec:background:moo}

A multi-objective optimization problem (MOP) seeks to simultaneously
optimize $m \geq 2$ objective functions
$f_1(x), \dots, f_m(x)$ over a decision space $x \in \Omega$. Because the
objectives typically conflict, there is generally no single $x$ that
minimizes (or maximizes) all of them at once; instead, a solution
$x_1$ is said to \emph{Pareto-dominate} $x_2$ if $x_1$ is at least as
good as $x_2$ on every objective and strictly better on at least one. A
solution is \emph{Pareto-optimal} if no other feasible solution
dominates it, and the image of the set of all Pareto-optimal solutions
in objective space is the \emph{Pareto front}. Rather than a single best
answer, a multi-objective evolutionary algorithm (MOEA) therefore aims to
return a population that approximates this front as closely and as
evenly as possible~\citep{emmerich2018tutorial}.

Comparing how well two algorithms' output populations approximate the
true (or best-known) Pareto front requires quality indicators that
capture both convergence and diversity, since neither alone is
sufficient: a tightly clustered set of points on the front converges well
but covers it poorly, while a widely spread set may cover a region no
solution actually dominates. The indicators used throughout this
survey are:
\begin{itemize}
  \item \textbf{Hypervolume (HV)}: the volume of objective space
    dominated by a solution set relative to a fixed reference point --
    the only indicator that is strictly monotonic with Pareto dominance
    on its own, and the basis of hypervolume-driven selection in
    SMS-EMOA~\citep{beume2007smsemoa}. Its value is sensitive to the
    reference point and normalization used, a methodological detail
    Section~\ref{sec:case-study:hypervolume} returns to.
  \item \textbf{Inverted Generational Distance (IGD)}: the average
    distance from each point on a reference (true or best-known) Pareto
    front to its nearest neighbor in the algorithm's output set,
    rewarding both convergence and coverage but requiring that reference
    front to be known or well-approximated -- a requirement this
    survey's own case study (Section~\ref{sec:case-study}) does not
    meet for any of its benchmark circuit families, so IGD does not
    itself recur there.
  \item \textbf{Spread and spacing}: measures of how evenly solutions are
    distributed along the approximated front, independent of how close
    that front is to the true optimum. Spread recurs in the case study
    (Section~\ref{sec:case-study:algorithms}); spacing does not.
\end{itemize}
See \citet{emmerich2018tutorial} and \citet{hanne2025review} for a fuller
treatment of these fundamentals and their formal definitions.

\subsection{Classical MOEAs}
\label{sec:background:moea}

The Pareto-dominance-based family of MOEAs -- the family every QIEA-MOO
method surveyed here either builds on or is benchmarked against -- traces
back to \citet{schaffer1985vega}'s Vector Evaluated Genetic Algorithm
(VEGA), which optimized each objective in a separate sub-population and
suffered from poor diversity as a result. Successive algorithms
addressed this diversity problem in turn: \citet{fonseca1993moga}'s MOGA
introduced Pareto-rank-based fitness with fitness sharing;
\citet{horn1994npga}'s NPGA added dominance-based tournament selection
with niching; \citet{zitzler1999spea}'s SPEA introduced an external
archive of non-dominated solutions and strength-based fitness;
\citet{knowles2000paes}'s PAES took a different design point entirely,
a $(1+1)$ evolution strategy with no population or crossover, using an
adaptive grid archive for diversity; and
\citet{corne2000pesa,corne2001pesa2}'s PESA and PESA-II introduced a
hyper-grid (``hyperbox'') archive with selection performed at the
hyperbox level. This lineage matters for QIEA-MOO specifically because
several surveyed methods graft a Q-bit/Q-gate representation onto one of
these selection mechanisms rather than inventing a new one --
understanding the mechanism being reused is a precondition for the
taxonomy's hybridization axis (Section~\ref{sec:taxonomy:hybridization}).

Five classical MOEAs recur most often as the baselines QIEA-MOO methods
are compared against in the surveyed literature, and this survey adopts
the same five as its background reference set:
\begin{itemize}
  \item \textbf{NSGA-II}~\citep{deb2002nsga2}: fast non-dominated
    sorting (ranking the population into successive Pareto fronts) combined
    with crowding-distance-based diversity preservation within each
    front; the most widely used MOEA baseline, and the one this survey's
    own case-study pipeline builds on (Section~\ref{sec:case-study}).
  \item \textbf{NSGA-III}~\citep{debjain2014nsga3}: replaces
    crowding distance with reference-point-based niching, extending
    NSGA-II's non-dominated-sorting framework to many-objective problems
    ($m \gg 3$) where crowding distance alone provides insufficient
    selection pressure.
  \item \textbf{SPEA2}~\citep{zitzler2001spea2}: refines SPEA with a
    finer-grained strength-based fitness assignment, a $k$-th-nearest-
    neighbor density estimator, and a fixed-size external archive.
  \item \textbf{MOEA/D}~\citep{zhangli2007moead}: decomposes the
    multi-objective problem into a set of scalar subproblems via weight
    vectors, solving them cooperatively using neighborhood information --
    a fundamentally different, decomposition-based design point from the
    three Pareto-dominance-based algorithms above.
  \item \textbf{SMS-EMOA}~\citep{beume2007smsemoa}: an indicator-based
    method whose environmental selection removes whichever individual
    contributes least to the population's dominated hypervolume, directly
    optimizing for the HV indicator described in
    Section~\ref{sec:background:moo} rather than a proxy for it.
\end{itemize}
Together these five span the main selection paradigms QIEA-MOO methods
adopt or compare against: crowding-distance-based, reference-point-based,
strength/density-based, decomposition-based, and hypervolume-driven.
Section~\ref{sec:taxonomy:selection} uses this same set of paradigms as
one of the taxonomy's classification axes.
